% \chapter{Bài toán đối ngẫu}

\section{Ý nghĩa và cách lập bài toán đối ngẫu}

\subsection{Ý nghĩa thực tế của bài toán đối ngẫu}
Trong thực tiễn kinh tế, mỗi bài toán Quy hoạch tuyến tính (gọi là bài toán gốc) luôn tồn tại song song một bài toán "anh em" với nó, được gọi là bài toán đối ngẫu. 

\begin{exampleblock}{Ví dụ minh họa}
Giả sử bạn là giám đốc của một công ty sản xuất (Công ty A) đang giải bài toán lập kế hoạch sản xuất tối ưu nhằm tối đa hóa lợi nhuận (tìm Max) dựa trên số lượng nguyên liệu giới hạn đang có trong kho [2]. Lúc này, một đối tác (Công ty B) đến đề nghị mua lại \textit{toàn bộ nguyên liệu} của bạn thay vì để bạn tự sản xuất [2].

Vấn đề đặt ra cho Công ty B là: \textit{Phải định giá mua mỗi đơn vị nguyên liệu là bao nhiêu để tổng chi phí mua là nhỏ nhất (tìm Min), nhưng vẫn đủ sức thuyết phục Công ty A bán lại (tức là giá mua nguyên liệu không được thấp hơn giá trị lợi nhuận mà Công ty A có thể tự làm ra)?} [2].

Bài toán định giá mua nguyên liệu của Công ty B chính là \textbf{Bài toán đối ngẫu} của bài toán sản xuất ban đầu [2].
\end{exampleblock}

\section{Quy tắc lập bài toán đối ngẫu}

Mối quan hệ chuyển đổi giữa bài toán gốc và bài toán đối ngẫu được thiết lập cực kỳ chặt chẽ. Một cách tổng quát, nếu bài toán gốc $(P)$ tìm giá trị lớn nhất (Max) thì bài toán đối ngẫu $(D)$ sẽ tìm giá trị nhỏ nhất (Min) và ngược lại [1]. 

Để thuận tiện cho việc lập mô hình, quy tắc chuyển đổi được tóm tắt qua hệ thống các bảng dưới đây:

\subsubsection{1. Bảng quy tắc tổng hợp bằng ký hiệu toán học}
Dưới đây là bảng ánh xạ 1-1 giữa các ràng buộc và biến số của hai bài toán được trình bày hoàn toàn bằng ngôn ngữ toán học [1].

\begin{table}[htbp]
    \centering
    \renewcommand{\arraystretch}{1.5}
    \begin{tabular}{|c|c|c|}
        \hline
        \textbf{Bài toán xuất phát (Gốc)} & $\Leftrightarrow$ & \textbf{Bài toán đối ngẫu} \\
        \hline
        \textbf{$Z = \sum_{j=1}^{n} c_j x_j \rightarrow \max$} & $\Leftrightarrow$ & \textbf{$Z' = \sum_{i=1}^{m} b_i y_i \rightarrow \min$} \\
        \hline
        $\sum_{j=1}^{n} a_{ij} x_j \le b_i$ & $\Leftrightarrow$ & $y_i \ge 0$ \\
        $\sum_{j=1}^{n} a_{ij} x_j \ge b_i$ & $\Leftrightarrow$ & $y_i \le 0$ \\
        $\sum_{j=1}^{n} a_{ij} x_j = b_i$ & $\Leftrightarrow$ & $y_i$ tùy ý \\
        \hline
        $x_j \ge 0$ & $\Leftrightarrow$ & $\sum_{i=1}^{m} a_{ij} y_i \ge c_j$ \\
        $x_j \le 0$ & $\Leftrightarrow$ & $\sum_{i=1}^{m} a_{ij} y_i \le c_j$ \\
        $x_j$ tùy ý & $\Leftrightarrow$ & $\sum_{i=1}^{m} a_{ij} y_i = c_j$ \\
        \hline
    \end{tabular}
    \caption{Bảng quy tắc chuyển đổi đối ngẫu (Dạng ký hiệu toán học)}
    \label{tab:doi_ngau_ky_hieu}
\end{table}

\textit{Mẹo ghi nhớ từ Bảng \ref{tab:doi_ngau_ky_hieu}:} Ràng buộc dấu (RBD) của bài toán này sẽ quyết định dấu của Ràng buộc chính (RBC) trong bài toán kia, và ngược lại [2].

\vspace{0.3cm}
\subsubsection{2. Phân tách quy tắc theo dạng hàm mục tiêu}
Để tránh nhầm lẫn trong quá trình làm bài, quy tắc trên được phân tách riêng biệt cho 2 dạng bài toán: từ Max sang Min và từ Min sang Max.

\textbf{Dạng 1: Chuyển từ bài toán Max sang Min} \\
Đối với bài toán gốc tìm Max, ràng buộc mang dấu ``$\le$'' và biến ``$\ge 0$'' được coi là các dấu "chuẩn".

\begin{table}[htbp]
    \centering
    \renewcommand{\arraystretch}{1.5}
    \begin{tabular}{|l|c|l|}
        \hline
        \multicolumn{1}{|c|}{\textbf{Bài toán Gốc (Max)}} & \textbf{$\rightarrow$} & \multicolumn{1}{c|}{\textbf{Bài toán Đối ngẫu (Min)}} \\
        \hline
        \textbf{Ràng buộc chính thứ $i$:} & & \textbf{Ràng buộc dấu của biến $y_i$:} \\
        1. Dấu ``$\le$'' (Dấu chuẩn) & $\rightarrow$ & 1. $y_i \ge 0$ (Dấu chuẩn) \\
        2. Dấu ``$\ge$'' (Ngược chuẩn) & $\rightarrow$ & 2. $y_i \le 0$ \\
        3. Dấu ``$=$'' & $\rightarrow$ & 3. $y_i$ tùy ý (không hạn chế dấu) \\
        \hline
        \textbf{Ràng buộc dấu của biến $x_j$:} & & \textbf{Ràng buộc chính thứ $j$:} \\
        1. $x_j \ge 0$ (Dấu chuẩn) & $\rightarrow$ & 1. Dấu ``$\ge$'' (Dấu chuẩn của Min) \\
        2. $x_j \le 0$ (Ngược chuẩn) & $\rightarrow$ & 2. Dấu ``$\le$'' \\
        3. $x_j$ tùy ý & $\rightarrow$ & 3. Dấu ``$=$'' \\
        \hline
    \end{tabular}
    \caption{Quy tắc chuyển đổi từ bài toán Max sang Min}
    \label{tab:doi_ngau_max_min}
\end{table}

\vspace{0.3cm}
\textbf{Dạng 2: Chuyển từ bài toán Min sang Max} \\
Đối với bài toán gốc tìm Min, mọi thứ đảo ngược lại: ràng buộc mang dấu ``$\ge$'' và biến ``$\ge 0$'' mới được coi là các dấu "chuẩn".

\begin{table}[htbp]
    \centering
    \renewcommand{\arraystretch}{1.5}
    \begin{tabular}{|l|c|l|}
        \hline
        \multicolumn{1}{|c|}{\textbf{Bài toán Gốc (Min)}} & \textbf{$\rightarrow$} & \multicolumn{1}{c|}{\textbf{Bài toán Đối ngẫu (Max)}} \\
        \hline
        \textbf{Ràng buộc chính thứ $i$:} & & \textbf{Ràng buộc dấu của biến $y_i$:} \\
        1. Dấu ``$\ge$'' (Dấu chuẩn) & $\rightarrow$ & 1. $y_i \ge 0$ (Dấu chuẩn) \\
        2. Dấu ``$\le$'' (Ngược chuẩn) & $\rightarrow$ & 2. $y_i \le 0$ \\
        3. Dấu ``$=$'' & $\rightarrow$ & 3. $y_i$ tùy ý (không hạn chế dấu) \\
        \hline
        \textbf{Ràng buộc dấu của biến $x_j$:} & & \textbf{Ràng buộc chính thứ $j$:} \\
        1. $x_j \ge 0$ (Dấu chuẩn) & $\rightarrow$ & 1. Dấu ``$\le$'' (Dấu chuẩn của Max) \\
        2. $x_j \le 0$ (Ngược chuẩn) & $\rightarrow$ & 2. Dấu ``$\ge$'' \\
        3. $x_j$ tùy ý & $\rightarrow$ & 3. Dấu ``$=$'' \\
        \hline
    \end{tabular}
    \caption{Quy tắc chuyển đổi từ bài toán Min sang Max}
    \label{tab:doi_ngau_min_max}
\end{table}

\subsection{Các định lý đối ngẫu cốt lõi}
Mối quan hệ chặt chẽ giữa bài toán gốc $(P)$ và bài toán đối ngẫu $(D)$ được thể hiện qua 3 định lý toán học cực kỳ quan trọng sau:

\begin{theorem}[\textbf{Định lý đối ngẫu yếu}]
Nếu $x^*$ là một phương án tùy ý của bài toán gốc $(P)$ tìm Max, và $y^*$ là một phương án tùy ý của bài toán đối ngẫu $(D)$ tìm Min, thì ta luôn có: $Z(x^*) \le Z'(y^*)$ [4]. \\
\textit{Hệ quả:} Nếu tìm được một cặp phương án $(x_0, y_0)$ sao cho $Z(x_0) = Z'(y_0)$, thì $x_0$ và $y_0$ chính là các phương án tối ưu của $(P)$ và $(D)$ [4].
\end{theorem}

\begin{theorem}[\textbf{Định lý đối ngẫu mạnh}]
Nếu bài toán gốc $(P)$ có phương án tối ưu $x^*$, thì bài toán đối ngẫu $(D)$ cũng chắc chắn có phương án tối ưu $y^*$, và giá trị tối ưu của hai hàm mục tiêu luôn bằng nhau: $Z_{max} = Z'_{min}$ [5].
\end{theorem}

\begin{theorem}[\textbf{Định lý độ lệch bù yếu (Complementary Slackness)}]
Điều kiện cần và đủ để phương án $x^0$ của $(P)$ và $y^0$ của $(D)$ cùng là Phương án tối ưu là chúng phải thỏa mãn hệ thức độ lệch bù. Nói một cách dễ hiểu, trong từng cặp ràng buộc đối ngẫu:
\begin{itemize}
    \item Nếu ràng buộc chính thứ $i$ của bài toán này \textbf{thỏa mãn lỏng} (dấu bất đẳng thức thực sự $>$ hoặc $<$), thì biến đối ngẫu tương ứng của bài toán kia bắt buộc phải \textbf{bằng 0} ($y_i = 0$ hoặc $x_i = 0$) [11].
    \item Nếu biến mang giá trị \textbf{khác 0} ($x_j > 0$ hoặc $y_i > 0$), thì ràng buộc chính tương ứng bên bài toán kia bắt buộc phải \textbf{thỏa mãn chặt} (xảy ra dấu bằng $=$) [11].
\end{itemize}
\end{theorem}

\section{Phương pháp Bảng cấu trúc ánh xạ}
    % \section{Phương pháp Bảng cấu trúc ánh xạ}

Để lập bài toán đối ngẫu một cách chính xác mà không bị nhầm lẫn về dấu, ta sử dụng phương pháp Bảng cấu trúc ánh xạ. Quá trình này giúp phân rã cấu trúc bài toán gốc và ánh xạ trực tiếp sang bài toán đối ngẫu.

\subsection{Quy tắc ánh xạ dấu}
Nguyên lý: \textit{Ràng buộc dấu (RBD) của bài toán này sẽ quyết định dấu của Ràng buộc chính (RBC) trong bài toán kia, và ngược lại.} Dựa vào dạng bài toán gốc, ta sẽ xác định được đâu là "dấu chuẩn" để khoanh tròn (\khoanh{}) và đâu là dấu ngược quy tắc để gạch bỏ (\cancel{}).

\textbf{1. Chuyển từ bài toán Max sang Min:}\\
Đối với bài toán gốc tìm Max, RBC mang dấu ``$\le$'' và biến ``$\ge 0$'' được coi là các dấu chuẩn.
\begin{itemize}
    \item \textbf{Từ RBC gốc $\rightarrow$ Dấu biến Đối ngẫu:} Dấu $\le \longrightarrow y_i \ge 0$ ; Dấu $\ge \longrightarrow y_i \le 0$ ; Dấu $=$ $\longrightarrow y_i$ tùy ý.
    \item \textbf{Từ Dấu biến gốc $\rightarrow$ RBC Đối ngẫu:} Biến $\ge 0 \longrightarrow$ Dấu $\ge$ ; Biến $\le 0 \longrightarrow$ Dấu $\le$ ; Biến tùy ý $\longrightarrow$ Dấu $=$.
\end{itemize}

\textbf{2. Chuyển từ bài toán Min sang Max:}\\
Đối với bài toán gốc tìm Min, RBC mang dấu ``$\ge$'' và biến ``$\ge 0$'' được coi là các dấu chuẩn.
\begin{itemize}
    \item \textbf{Từ RBC gốc $\rightarrow$ Dấu biến Đối ngẫu:} Dấu $\ge \longrightarrow y_i \ge 0$ ; Dấu $\le \longrightarrow y_i \le 0$ ; Dấu $=$ $\longrightarrow y_i$ tùy ý.
    \item \textbf{Từ Dấu biến gốc $\rightarrow$ RBC Đối ngẫu:} Biến $\ge 0 \longrightarrow$ Dấu $\le$ ; Biến $\le 0 \longrightarrow$ Dấu $\ge$ ; Biến tùy ý $\longrightarrow$ Dấu $=$.
\end{itemize}

\subsection{Biểu mẫu cấu trúc bảng ánh xạ}
Dưới đây là biểu mẫu ma trận ánh xạ tổng quát cho bài toán $3 \times 3$. Trong thực hành, ta lập Bảng phân rã cho bài toán Gốc, sau đó bốc các hệ số dựa vào ký hiệu để điền vào Bảng cấu trúc Đối ngẫu và khoanh dấu.

\textbf{Bước 1: Lập bảng phân rã cấu trúc bài toán Gốc (P)}
\begin{table}[H]
    \centering
    \renewcommand{\arraystretch}{1.5}
    \resizebox{\textwidth}{!}{
    \begin{tabular}{|c|c c c c c c c c|c c c|c|}
        \hline
        \textbf{Hàm/RBC} & \multicolumn{8}{c|}{\textbf{Hệ số \& Biến số}} & \multicolumn{3}{c|}{\textbf{Dấu RBC}} & \textbf{Hệ số $b$} \\
        \hline
        $Z (\max/\min)$ & \TL{c_1}{} & $x_1$ & $+$ & \TL{c_2}{} & $x_2$ & $+$ & \TL{c_3}{} & $x_3$ & \multicolumn{3}{c|}{} & \\
        \hline
        $R_1$ & \TL{a_{11}}{} & $x_1$ & $+$ & \TL{a_{12}}{} & $x_2$ & $+$ & \TL{a_{13}}{} & $x_3$ & $\ge$ & $\le$ & $=$ & \TL{b_1}{} \\
        \hline
        $R_2$ & \TL{a_{21}}{} & $x_1$ & $+$ & \TL{a_{22}}{} & $x_2$ & $+$ & \TL{a_{23}}{} & $x_3$ & $\ge$ & $\le$ & $=$ & \TL{b_2}{} \\
        \hline
        $R_3$ & \TL{a_{31}}{} & $x_1$ & $+$ & \TL{a_{32}}{} & $x_2$ & $+$ & \TL{a_{33}}{} & $x_3$ & $\ge$ & $\le$ & $=$ & \TL{b_3}{} \\
        \hline
        \multirow{3}{*}{\shortstack{Dấu \\ Biến số}} 
        & \multicolumn{2}{c}{$x_1$} & \multicolumn{6}{l|}{$\ge 0$ \quad $\le 0$ \quad Tùy ý} & \multicolumn{4}{c|}{} \\
        & \multicolumn{2}{c}{$x_2$} & \multicolumn{6}{l|}{$\ge 0$ \quad $\le 0$ \quad Tùy ý} & \multicolumn{4}{c|}{} \\
        & \multicolumn{2}{c}{$x_3$} & \multicolumn{6}{l|}{$\ge 0$ \quad $\le 0$ \quad Tùy ý} & \multicolumn{4}{c|}{} \\
        \hline
    \end{tabular}
    }
    \caption{Bảng phân rã cấu trúc bài toán Gốc (P)}
\end{table}

\textbf{Bước 2: Lập bảng cấu trúc bài toán Đối ngẫu (D)} \\
\textit{(Áp dụng quy tắc ánh xạ ở mục trên để gạch bỏ $\cancel{}$ hoặc khoanh tròn \khoanh{} các dấu tương ứng)}

\begin{table}[H]
    \centering
    \renewcommand{\arraystretch}{1.5}
    \resizebox{\textwidth}{!}{
    \begin{tabular}{|c|c c c c c c c c|c c c|c|}
        \hline
        \textbf{Hàm/RBC} & \multicolumn{8}{c|}{\textbf{Hệ số \& Biến số}} & \multicolumn{3}{c|}{\textbf{Dấu RBC ĐN}} & \textbf{Hệ số $c$} \\
        \hline
        $Z' (\min/\max)$ & \TL{b_1}{} & $y_1$ & $+$ & \TL{b_2}{} & $y_2$ & $+$ & \TL{b_3}{} & $y_3$ & \multicolumn{3}{c|}{} & \\
        \hline
        $R'_1$ & \TL{a_{11}}{} & $y_1$ & $+$ & \TL{a_{21}}{} & $y_2$ & $+$ & \TL{a_{31}}{} & $y_3$ & $\ge$ & $\le$ & $=$ & \TL{c_1}{} \\
        \hline
        $R'_2$ & \TL{a_{12}}{} & $y_1$ & $+$ & \TL{a_{22}}{} & $y_2$ & $+$ & \TL{a_{32}}{} & $y_3$ & $\ge$ & $\le$ & $=$ & \TL{c_2}{} \\
        \hline
        $R'_3$ & \TL{a_{13}}{} & $y_1$ & $+$ & \TL{a_{23}}{} & $y_2$ & $+$ & \TL{a_{33}}{} & $y_3$ & $\ge$ & $\le$ & $=$ & \TL{c_3}{} \\
        \hline
        \multirow{3}{*}{\shortstack{Dấu Biến \\ Đối ngẫu}} 
        & \multicolumn{2}{c}{$y_1$} & \multicolumn{6}{l|}{$\ge 0$ \quad $\le 0$ \quad Tùy ý} & \multicolumn{4}{c|}{} \\
        & \multicolumn{2}{c}{$y_2$} & \multicolumn{6}{l|}{$\ge 0$ \quad $\le 0$ \quad Tùy ý} & \multicolumn{4}{c|}{} \\
        & \multicolumn{2}{c}{$y_3$} & \multicolumn{6}{l|}{$\ge 0$ \quad $\le 0$ \quad Tùy ý} & \multicolumn{4}{c|}{} \\
        \hline
    \end{tabular}
    }
    \caption{Bảng cấu trúc bài toán Đối ngẫu (D) sau khi ánh xạ}
\end{table}

\section{Phương pháp Bảng cấu trúc ánh xạ trực tiếp}

\subsection{Biểu mẫu ánh xạ: Bài toán gốc (MAX) \texorpdfstring{$\longrightarrow$}{→} Đối ngẫu (MIN)}

Đối với bài toán gốc MAX, dấu chuẩn của Ràng buộc chính (RBC) là $\le$ (sinh ra biến đối ngẫu $\ge 0$) và dấu chuẩn của biến là $\ge 0$ (sinh ra RBC đối ngẫu $\ge$).

\begin{table}[H]
    \centering
    \renewcommand{\arraystretch}{1.8}
    \resizebox{\textwidth}{!}{
    \begin{tabular}{|c|c c c c c c c c|c|c|}
        \hline
        \textbf{Bài toán GỐC} & \multicolumn{8}{c|}{\textbf{Hệ số \& Biến số gốc ($x_j$)}} & \textbf{Ánh xạ Dấu} & \textbf{Hệ số $b$} \\
        \hline
        $Z (\max)$ & \kh{blue}{$c_1$} & $x_1$ & $+$ & \kh{blue}{$c_2$} & $x_2$ & $+$ & \kh{blue}{$c_3$} & $x_3$ & \multicolumn{2}{c|}{\textit{(Đối ngẫu: $Z' \rightarrow \min$)}} \\
        \hline
        \textbf{RBC $\rightarrow$ Biến ĐN} & \multicolumn{8}{c|}{\textit{Ma trận hệ số $A$ chuyển vị thành $A^T$}} & \textbf{RBC gốc $\rightarrow$ Dấu biến ĐN ($y_i$)} & \\
        \hline
        $R_1 \longrightarrow y_1$ & \khdouble{green}{$a_{11}$} & $x_1$ & $+$ & \khdouble{green}{$a_{12}$} & $x_2$ & $+$ & \khdouble{green}{$a_{13}$} & $x_3$ & 
        $\le \longrightarrow y_1 \ge 0$ \quad $\ge \longrightarrow y_1 \le 0$ \quad $= \longrightarrow$ tùy ý & \khdash{red}{$b_1$} \\
        \hline
        $R_2 \longrightarrow y_2$ & $a_{21}$ & $x_1$ & $+$ & $a_{22}$ & $x_2$ & $+$ & $a_{23}$ & $x_3$ & 
        $\le \longrightarrow y_2 \ge 0$ \quad $\ge \longrightarrow y_2 \le 0$ \quad $= \longrightarrow$ tùy ý & \khdash{red}{$b_2$} \\
        \hline
        $R_3 \longrightarrow y_3$ & $a_{31}$ & $x_1$ & $+$ & $a_{32}$ & $x_2$ & $+$ & $a_{33}$ & $x_3$ & 
        $\le \longrightarrow y_3 \ge 0$ \quad $\ge \longrightarrow y_3 \le 0$ \quad $= \longrightarrow$ tùy ý & \khdash{red}{$b_3$} \\
        \hline
        \multirow{4}{*}{\shortstack{\textbf{Dấu Biến gốc} \\ $\downarrow$ \\ \textbf{Dấu RBC ĐN}}} 
        & \multicolumn{2}{c}{$x_1$} & \multicolumn{6}{l|}{$\ge 0 \longrightarrow \ge$ \qquad $\le 0 \longrightarrow \le$ \qquad Tùy ý $\longrightarrow =$ } & \multicolumn{2}{c|}{ \multirow{4}{*}{\shortstack{\textit{Các cột này} \\ \textit{trở thành} \\ \textit{RBC đối ngẫu} \\ \textit{$R'_1, R'_2, R'_3$}}} } \\
        \cline{2-9}
        & \multicolumn{2}{c}{$x_2$} & \multicolumn{6}{l|}{$\ge 0 \longrightarrow \ge$ \qquad $\le 0 \longrightarrow \le$ \qquad Tùy ý $\longrightarrow =$ } & \multicolumn{2}{c|}{} \\
        \cline{2-9}
        & \multicolumn{2}{c}{$x_3$} & \multicolumn{6}{l|}{$\ge 0 \longrightarrow \ge$ \qquad $\le 0 \longrightarrow \le$ \qquad Tùy ý $\longrightarrow =$ } & \multicolumn{2}{c|}{} \\
        \hline
    \end{tabular}
    }
    \caption{Bảng phân rã và ánh xạ trực tiếp cho bài toán gốc MAX}
\end{table}

\begin{table}[H]
    \centering
    \renewcommand{\arraystretch}{1.8}
    \resizebox{\textwidth}{!}{
    \begin{tabular}{|c|c c c c c c c c|c|c|}
        \hline
        \textbf{Bài toán ĐỐI NGẪU} & \multicolumn{8}{c|}{\textbf{Hệ số \& Biến số ĐN ($y_i$)}} & \textbf{Dấu RBC ĐN} & \textbf{Hệ số $c$} \\
        \hline
        $Z' (\min)$ & \khdash{red}{$b_1$} & $y_1$ & $+$ & \khdash{red}{$b_2$} & $y_2$ & $+$ & \khdash{red}{$b_3$} & $y_3$ & \multicolumn{2}{c|}{\textit{(Nguồn: Hệ số tự do gốc)}} \\
        \hline
        \textbf{Biến gốc $\rightarrow$ RBC ĐN} & \multicolumn{8}{c|}{\textit{Ma trận hệ số chuyển vị $A^T$}} & \textbf{Dấu biến gốc $\rightarrow$ Dấu RBC ĐN} & \\
        \hline
        $x_1 \longrightarrow R'_1$ & \khdouble{green}{$a_{11}$} & $y_1$ & $+$ & $a_{21}$ & $y_2$ & $+$ & $a_{31}$ & $y_3$ & 
        $\ge 0 \longrightarrow \ge$ \quad $\le 0 \longrightarrow \le$ \quad Tùy ý $\longrightarrow =$ & \kh{blue}{$c_1$} \\
        \hline
        $x_2 \longrightarrow R'_2$ & \khdouble{green}{$a_{12}$} & $y_1$ & $+$ & $a_{22}$ & $y_2$ & $+$ & $a_{32}$ & $y_3$ & 
        $\ge 0 \longrightarrow \ge$ \quad $\le 0 \longrightarrow \le$ \quad Tùy ý $\longrightarrow =$ & \kh{blue}{$c_2$} \\
        \hline
        $x_3 \longrightarrow R'_3$ & \khdouble{green}{$a_{13}$} & $y_1$ & $+$ & $a_{23}$ & $y_2$ & $+$ & $a_{33}$ & $y_3$ & 
        $\ge 0 \longrightarrow \ge$ \quad $\le 0 \longrightarrow \le$ \quad Tùy ý $\longrightarrow =$ & \kh{blue}{$c_3$} \\
        \hline
        \multirow{4}{*}{\shortstack{\textbf{Dấu RBC gốc} \\ $\downarrow$ \\ \textbf{Dấu Biến ĐN}}} 
        & \multicolumn{2}{c}{$y_1$} & \multicolumn{6}{l|}{$\le \longrightarrow \ge 0$ \qquad $\ge \longrightarrow \le 0$ \qquad $= \longrightarrow$ Tùy ý } & \multicolumn{2}{c|}{ \multirow{4}{*}{\shortstack{\textit{Ma trận A} \\ \textit{đã được} \\ \textit{xoay ngang} \\ \textit{thành dọc}}} } \\
        \cline{2-9}
        & \multicolumn{2}{c}{$y_2$} & \multicolumn{6}{l|}{$\le \longrightarrow \ge 0$ \qquad $\ge \longrightarrow \le 0$ \qquad $= \longrightarrow$ Tùy ý } & \multicolumn{2}{c|}{} \\
        \cline{2-9}
        & \multicolumn{2}{c}{$y_3$} & \multicolumn{6}{l|}{$\le \longrightarrow \ge 0$ \qquad $\ge \longrightarrow \le 0$ \qquad $= \longrightarrow$ Tùy ý } & \multicolumn{2}{c|}{} \\
        \hline
    \end{tabular}
    }
    \caption{Bảng phân rã bài toán Đối ngẫu (MIN) và ánh xạ hệ số từ bài toán Gốc}
\end{table}

\subsubsection*{Ví dụ áp dụng: Bài toán lập kế hoạch sản xuất}
Xét bài toán gốc (P) tìm phương án sản xuất tối ưu để đạt lợi nhuận lớn nhất:
\begin{align*}
    Z &= 11x_1 + 9x_2 + 15x_3 \rightarrow \max \\
\text{sao cho:} \quad
    &5x_1 + 4x_2 + 4x_3 \le 21000 \quad (R_1)\\
    &3x_1 + 6x_2 + 3x_3 \le 15000 \quad (R_2)\\
    &3x_1 + 7x_2 + 5x_3 \le 26000 \quad (R_3)\\
    &x_1, x_2, x_3 \ge 0
\end{align*}

Dưới đây là quá trình phân rã dữ kiện và ánh xạ hệ số từ bài toán Gốc (P) sang bài toán Đối ngẫu (D):

\begin{table}[H]
    \centering
    \renewcommand{\arraystretch}{1.8}
    \resizebox{\textwidth}{!}{
    \begin{tabular}{|c|c c c c c c c c|c|c|}
        \hline
        \textbf{Bài toán GỐC} & \multicolumn{8}{c|}{\textbf{Hệ số \& Biến số gốc ($x_j$)}} & \textbf{Ánh xạ Dấu} & \textbf{Hệ số $b$} \\
        \hline
        $Z (\max)$ & \TL{c_1}{\kh{blue}{11}} & $x_1$ & $+$ & \TL{c_2}{\kh{blue}{9}} & $x_2$ & $+$ & \TL{c_3}{\kh{blue}{15}} & $x_3$ & \multicolumn{2}{c|}{\textit{(Đối ngẫu: $Z' \rightarrow \min$)}} \\
        \hline
        \textbf{RBC $\rightarrow$ Biến ĐN} & \multicolumn{8}{c|}{\textit{Ma trận hệ số $A$ chuyển vị thành $A^T$}} & \textbf{RBC gốc $\rightarrow$ Dấu biến ĐN ($y_i$)} & \\
        \hline
        $R_1 \longrightarrow y_1$ & \TL{a_{11}}{\khdash{red}{5}} & $x_1$ & $+$ & \TL{a_{12}}{\khdash{red}{4}} & $x_2$ & $+$ & \TL{a_{13}}{\khdash{red}{4}} & $x_3$ & 
        \kv{green}{$\le$} / $\ge$ / $=$ $\longrightarrow$ \kv{green}{$\ge 0$} / $\le 0$ / Tùy ý & \TL{b_1}{\koval{green}{21000}} \\
        \hline
        $R_2 \longrightarrow y_2$ & \TL{a_{21}}{3} & $x_1$ & $+$ & \TL{a_{22}}{6} & $x_2$ & $+$ & \TL{a_{23}}{3} & $x_3$ & 
        \kv{green}{$\le$} / $\ge$ / $=$ $\longrightarrow$ \kv{green}{$\ge 0$} / $\le 0$ / Tùy ý & \TL{b_2}{\koval{green}{15000}} \\
        \hline
        $R_3 \longrightarrow y_3$ & \TL{a_{31}}{3} & $x_1$ & $+$ & \TL{a_{32}}{7} & $x_2$ & $+$ & \TL{a_{33}}{5} & $x_3$ & 
        \kv{green}{$\le$} / $\ge$ / $=$ $\longrightarrow$ \kv{green}{$\ge 0$} / $\le 0$ / Tùy ý & \TL{b_3}{\koval{green}{26000}} \\
        \hline
        \multirow{4}{*}{\shortstack{\textbf{Dấu Biến gốc} \\ $\downarrow$ \\ \textbf{Dấu RBC ĐN}}} 
        & \multicolumn{2}{c}{$x_1$} & \multicolumn{6}{l|}{\kvd{red}{$\ge 0$} / $\le 0$ / Tùy ý $\longrightarrow$ \kvd{red}{$\ge$} / $\le$ / $=$ } & \multicolumn{2}{c|}{ \multirow{4}{*}{\shortstack{\textit{Các cột này} \\ \textit{trở thành} \\ \textit{RBC đối ngẫu} \\ \textit{$R'_1, R'_2, R'_3$}}} } \\
        \cline{2-9}
        & \multicolumn{2}{c}{$x_2$} & \multicolumn{6}{l|}{\kvd{red}{$\ge 0$} / $\le 0$ / Tùy ý $\longrightarrow$ \kvd{red}{$\ge$} / $\le$ / $=$ } & \multicolumn{2}{c|}{} \\
        \cline{2-9}
        & \multicolumn{2}{c}{$x_3$} & \multicolumn{6}{l|}{\kvd{red}{$\ge 0$} / $\le 0$ / Tùy ý $\longrightarrow$ \kvd{red}{$\ge$} / $\le$ / $=$ } & \multicolumn{2}{c|}{} \\
        \hline
    \end{tabular}
    }
    \caption{Phân rã dữ kiện và cấu trúc dấu của bài toán Gốc (MAX)}
\end{table}

\begin{table}[H]
    \centering
    \renewcommand{\arraystretch}{1.8}
    \resizebox{\textwidth}{!}{
    \begin{tabular}{|c|c c c c c c c c|c|c|}
        \hline
        \textbf{Bài toán ĐỐI NGẪU} & \multicolumn{8}{c|}{\textbf{Hệ số \& Biến số ĐN ($y_i$)}} & \textbf{Dấu RBC ĐN} & \textbf{Hệ số $c$} \\
        \hline
        $Z' (\min)$ & \TL{b_1}{\koval{green}{21000}} & $y_1$ & $+$ & \TL{b_2}{\koval{green}{15000}} & $y_2$ & $+$ & \TL{b_3}{\koval{green}{26000}} & $y_3$ & \multicolumn{2}{c|}{\textit{(Từ hệ số mục tiêu gốc)}} \\
        \hline
        \textbf{Biến gốc $\rightarrow$ RBC ĐN} & \multicolumn{8}{c|}{\textit{Ma trận hệ số chuyển vị $A^T$}} & \textbf{Dấu biến gốc $\rightarrow$ Dấu RBC ĐN} & \\
        \hline
        $x_1 \longrightarrow R'_1$ & \TL{a_{11}}{\khdash{red}{5}} & $y_1$ & $+$ & \TL{a_{21}}{3} & $y_2$ & $+$ & \TL{a_{31}}{3} & $y_3$ & 
        \kvd{red}{$\ge 0$} / $\le 0$ / Tùy ý $\longrightarrow$ \kvd{red}{$\ge$} / $\le$ / $=$ & \TL{c_1}{\kh{blue}{11}} \\
        \hline
        $x_2 \longrightarrow R'_2$ & \TL{a_{12}}{\khdash{red}{4}} & $y_1$ & $+$ & \TL{a_{22}}{6} & $y_2$ & $+$ & \TL{a_{32}}{7} & $y_3$ & 
        \kvd{red}{$\ge 0$} / $\le 0$ / Tùy ý $\longrightarrow$ \kvd{red}{$\ge$} / $\le$ / $=$ & \TL{c_2}{\kh{blue}{9}} \\
        \hline
        $x_3 \longrightarrow R'_3$ & \TL{a_{13}}{\khdash{red}{4}} & $y_1$ & $+$ & \TL{a_{23}}{3} & $y_2$ & $+$ & \TL{a_{33}}{5} & $y_3$ & 
        \kvd{red}{$\ge 0$} / $\le 0$ / Tùy ý $\longrightarrow$ \kvd{red}{$\ge$} / $\le$ / $=$ & \TL{c_3}{\kh{blue}{15}} \\
        \hline
        \multirow{4}{*}{\shortstack{\textbf{Dấu RBC gốc} \\ $\downarrow$ \\ \textbf{Dấu Biến ĐN}}} 
        & \multicolumn{2}{c}{$y_1$} & \multicolumn{6}{l|}{\kv{green}{$\le$} / $\ge$ / $=$ $\longrightarrow$ \kv{green}{$\ge 0$} / $\le 0$ / Tùy ý } & \multicolumn{2}{c|}{ \multirow{4}{*}{\shortstack{\textit{Các cột này} \\ \textit{thể hiện} \\ \textit{chiều biến} \\ \textit{đối ngẫu}}} } \\
        \cline{2-9}
        & \multicolumn{2}{c}{$y_2$} & \multicolumn{6}{l|}{\kv{green}{$\le$} / $\ge$ / $=$ $\longrightarrow$ \kv{green}{$\ge 0$} / $\le 0$ / Tùy ý } & \multicolumn{2}{c|}{} \\
        \cline{2-9}
        & \multicolumn{2}{c}{$y_3$} & \multicolumn{6}{l|}{\kv{green}{$\le$} / $\ge$ / $=$ $\longrightarrow$ \kv{green}{$\ge 0$} / $\le 0$ / Tùy ý } & \multicolumn{2}{c|}{} \\
        \hline
    \end{tabular}
    }
    \caption{Ánh xạ hệ số và thiết lập bài toán Đối ngẫu (MIN)}
\end{table}

\noindent Từ Bảng 2, ta dễ dàng viết được bài toán Đối ngẫu (D) một cách trực quan và chính xác:
\begin{align*}
    Z' &= 21000y_1 + 15000y_2 + 26000y_3 \rightarrow \min \\
\text{sao cho:} \quad
    &5y_1 + 3y_2 + 3y_3 \ge 11 \\
    &4y_1 + 6y_2 + 7y_3 \ge 9 \\
    &4y_1 + 3y_2 + 5y_3 \ge 15 \\
    &y_1, y_2, y_3 \ge 0
\end{align*}
\subsection{Biểu mẫu ánh xạ: Bài toán gốc (MIN) \texorpdfstring{$\longrightarrow$}{→} Đối ngẫu (MAX)}

Đối với bài toán gốc MIN, mọi quy tắc đảo ngược lại: dấu chuẩn của RBC là $\ge$ (sinh ra biến đối ngẫu $\ge 0$) và dấu chuẩn của biến là $\ge 0$ (sinh ra RBC đối ngẫu $\le$).

\begin{table}[H]
    \centering
    \renewcommand{\arraystretch}{1.8}
    \resizebox{\textwidth}{!}{
    \begin{tabular}{|c|c c c c c c c c|c|c|}
        \hline
        \textbf{Bài toán GỐC} & \multicolumn{8}{c|}{\textbf{Hệ số \& Biến số gốc ($x_j$)}} & \textbf{Ánh xạ Dấu} & \textbf{Hệ số $b$} \\
        \hline
        $Z (\min)$ & \kh{blue}{$c_1$} & $x_1$ & $+$ & \kh{blue}{$c_2$} & $x_2$ & $+$ & \kh{blue}{$c_3$} & $x_3$ & \multicolumn{2}{c|}{\textit{(Đối ngẫu: $Z' \rightarrow \max$)}} \\
        \hline
        \textbf{RBC $\rightarrow$ Biến ĐN} & \multicolumn{8}{c|}{\textit{Ma trận hệ số $A$ chuyển vị thành $A^T$}} & \textbf{RBC gốc $\rightarrow$ Dấu biến ĐN ($y_i$)} & \\
        \hline
        $R_1 \longrightarrow y_1$ & \kh{green}{$a_{11}$} & $x_1$ & $+$ & \kh{green}{$a_{12}$} & $x_2$ & $+$ & \kh{green}{$a_{13}$} & $x_3$ & 
        $\ge \longrightarrow y_1 \ge 0$ \quad $\le \longrightarrow y_1 \le 0$ \quad $= \longrightarrow$ tùy ý & \kh{red}{$b_1$} \\
        \hline
        $R_2 \longrightarrow y_2$ & \kh{green}{$a_{21}$} & $x_1$ & $+$ & \kh{green}{$a_{22}$} & $x_2$ & $+$ & \kh{green}{$a_{23}$} & $x_3$ & 
        $\ge \longrightarrow y_2 \ge 0$ \quad $\le \longrightarrow y_2 \le 0$ \quad $= \longrightarrow$ tùy ý & \kh{red}{$b_2$} \\
        \hline
        $R_3 \longrightarrow y_3$ & \kh{green}{$a_{31}$} & $x_1$ & $+$ & \kh{green}{$a_{32}$} & $x_2$ & $+$ & \kh{green}{$a_{33}$} & $x_3$ & 
        $\ge \longrightarrow y_3 \ge 0$ \quad $\le \longrightarrow y_3 \le 0$ \quad $= \longrightarrow$ tùy ý & \kh{red}{$b_3$} \\
        \hline
        \multirow{4}{*}{\shortstack{\textbf{Dấu Biến gốc} \\ $\downarrow$ \\ \textbf{Dấu RBC ĐN}}} 
        & \multicolumn{2}{c}{$x_1$} & \multicolumn{6}{l|}{$\ge 0 \longrightarrow \le$ \qquad $\le 0 \longrightarrow \ge$ \qquad Tùy ý $\longrightarrow =$ } & \multicolumn{2}{c|}{ \multirow{4}{*}{\shortstack{\textit{Các cột này} \\ \textit{trở thành} \\ \textit{RBC đối ngẫu} \\ \textit{$R'_1, R'_2, R'_3$}}} } \\
        \cline{2-9}
        & \multicolumn{2}{c}{$x_2$} & \multicolumn{6}{l|}{$\ge 0 \longrightarrow \le$ \qquad $\le 0 \longrightarrow \ge$ \qquad Tùy ý $\longrightarrow =$ } & \multicolumn{2}{c|}{} \\
        \cline{2-9}
        & \multicolumn{2}{c}{$x_3$} & \multicolumn{6}{l|}{$\ge 0 \longrightarrow \le$ \qquad $\le 0 \longrightarrow \ge$ \qquad Tùy ý $\longrightarrow =$ } & \multicolumn{2}{c|}{} \\
        \hline
    \end{tabular}
    }
    \caption{Bảng phân rã và ánh xạ trực tiếp cho bài toán gốc MIN}
\end{table}

\begin{table}[H]
    \centering
    \renewcommand{\arraystretch}{1.8}
    \resizebox{\textwidth}{!}{
    \begin{tabular}{|c|c c c c c c c c|c|c|}
        \hline
        \textbf{Bài toán ĐỐI NGẪU} & \multicolumn{8}{c|}{\textbf{Hệ số \& Biến số ĐN ($y_i$)}} & \textbf{Dấu RBC ĐN} & \textbf{Hệ số $c$} \\
        \hline
        $Z' (\max)$ & \kh{red}{$b_1$} & $y_1$ & $+$ & \kh{red}{$b_2$} & $y_2$ & $+$ & \kh{red}{$b_3$} & $y_3$ & \multicolumn{2}{c|}{\textit{(Nguồn: Hệ số tự do gốc)}} \\
        \hline
        \textbf{Biến gốc $\rightarrow$ RBC ĐN} & \multicolumn{8}{c|}{\textit{Ma trận hệ số chuyển vị $A^T$}} & \textbf{Dấu biến gốc $\rightarrow$ Dấu RBC ĐN} & \\
        \hline
        $x_1 \longrightarrow R'_1$ & \kh{green}{$a_{11}$} & $y_1$ & $+$ & \kh{green}{$a_{21}$} & $y_2$ & $+$ & \kh{green}{$a_{31}$} & $y_3$ & 
        $\ge 0 \longrightarrow \le$ \quad $\le 0 \longrightarrow \ge$ \quad Tùy ý $\longrightarrow =$ & \kh{blue}{$c_1$} \\
        \hline
        $x_2 \longrightarrow R'_2$ & \kh{green}{$a_{12}$} & $y_1$ & $+$ & \kh{green}{$a_{22}$} & $y_2$ & $+$ & \kh{green}{$a_{32}$} & $y_3$ & 
        $\ge 0 \longrightarrow \le$ \quad $\le 0 \longrightarrow \ge$ \quad Tùy ý $\longrightarrow =$ & \kh{blue}{$c_2$} \\
        \hline
        $x_3 \longrightarrow R'_3$ & \kh{green}{$a_{13}$} & $y_1$ & $+$ & \kh{green}{$a_{23}$} & $y_2$ & $+$ & \kh{green}{$a_{33}$} & $y_3$ & 
        $\ge 0 \longrightarrow \le$ \quad $\le 0 \longrightarrow \ge$ \quad Tùy ý $\longrightarrow =$ & \kh{blue}{$c_3$} \\
        \hline
        \multirow{4}{*}{\shortstack{\textbf{Dấu RBC gốc} \\ $\downarrow$ \\ \textbf{Dấu Biến ĐN}}} 
        & \multicolumn{2}{c}{$y_1$} & \multicolumn{6}{l|}{$\ge \longrightarrow \ge 0$ \qquad $\le \longrightarrow \le 0$ \qquad $= \longrightarrow$ Tùy ý } & \multicolumn{2}{c|}{ \multirow{4}{*}{\shortstack{\textit{Ma trận A} \\ \textit{đã được} \\ \textit{xoay ngang} \\ \textit{thành dọc}}} } \\
        \cline{2-9}
        & \multicolumn{2}{c}{$y_2$} & \multicolumn{6}{l|}{$\ge \longrightarrow \ge 0$ \qquad $\le \longrightarrow \le 0$ \qquad $= \longrightarrow$ Tùy ý } & \multicolumn{2}{c|}{} \\
        \cline{2-9}
        & \multicolumn{2}{c}{$y_3$} & \multicolumn{6}{l|}{$\ge \longrightarrow \ge 0$ \qquad $\le \longrightarrow \le 0$ \qquad $= \longrightarrow$ Tùy ý } & \multicolumn{2}{c|}{} \\
        \hline
    \end{tabular}
    }
    \caption{Bảng phân rã bài toán Đối ngẫu (MAX) và ánh xạ hệ số từ bài toán Gốc}
\end{table}

\subsubsection*{Ví dụ áp dụng: Bài toán chi phí tối thiểu}
Xét bài toán gốc (P) tìm phương án tối ưu để đạt chi phí nhỏ nhất:
\begin{align*}
    Z &= 12x_1 + 27x_2 + 6x_3 \rightarrow \min \\
\text{sao cho:} \quad
    &2x_1 + 3x_2 + 2x_3 \ge 12 \quad (R_1)\\
    &x_1 + 3x_2 + x_3 \ge 6 \quad (R_2)\\
    &6x_1 + 9x_2 + 2x_3 \ge 24 \quad (R_3)\\
    &x_1, x_2, x_3 \ge 0
\end{align*}

Dưới đây là quá trình phân rã dữ kiện và ánh xạ hệ số từ bài toán Gốc (P) sang bài toán Đối ngẫu (D):

\begin{table}[H]
    \centering
    \renewcommand{\arraystretch}{1.8}
    \resizebox{\textwidth}{!}{
    \begin{tabular}{|c|c c c c c c c c|c|c|}
        \hline
        \textbf{Bài toán GỐC} & \multicolumn{8}{c|}{\textbf{Hệ số \& Biến số gốc ($x_j$)}} & \textbf{Ánh xạ Dấu} & \textbf{Hệ số $b$} \\
        \hline
        $Z (\min)$ & \TL{c_1}{12} & $x_1$ & $+$ & \TL{c_2}{27} & $x_2$ & $+$ & \TL{c_3}{6} & $x_3$ & \multicolumn{2}{c|}{\textit{(Đối ngẫu: $Z' \rightarrow \max$)}} \\
        \hline
        \textbf{RBC $\rightarrow$ Biến ĐN} & \multicolumn{8}{c|}{\textit{Ma trận hệ số $A$ chuyển vị thành $A^T$}} & \textbf{RBC gốc $\rightarrow$ Dấu biến ĐN ($y_i$)} & \\
        \hline
        $R_1 \longrightarrow y_1$ & \TL{a_{11}}{2} & $x_1$ & $+$ & \TL{a_{12}}{3} & $x_2$ & $+$ & \TL{a_{13}}{2} & $x_3$ & 
        \khoanh{$\ge$} / $\le$ / $=$ $\longrightarrow$ \khoanh{$\ge 0$} / $\le 0$ / Tùy ý & \TL{b_1}{12} \\
        \hline
        $R_2 \longrightarrow y_2$ & \TL{a_{21}}{1} & $x_1$ & $+$ & \TL{a_{22}}{3} & $x_2$ & $+$ & \TL{a_{23}}{1} & $x_3$ & 
        \khoanh{$\ge$} / $\le$ / $=$ $\longrightarrow$ \khoanh{$\ge 0$} / $\le 0$ / Tùy ý & \TL{b_2}{6} \\
        \hline
        $R_3 \longrightarrow y_3$ & \TL{a_{31}}{6} & $x_1$ & $+$ & \TL{a_{32}}{9} & $x_2$ & $+$ & \TL{a_{33}}{2} & $x_3$ & 
        \khoanh{$\ge$} / $\le$ / $=$ $\longrightarrow$ \khoanh{$\ge 0$} / $\le 0$ / Tùy ý & \TL{b_3}{24} \\
        \hline
        \multirow{4}{*}{\shortstack{\textbf{Dấu Biến gốc} \\ $\downarrow$ \\ \textbf{Dấu RBC ĐN}}} 
        & \multicolumn{2}{c}{$x_1$} & \multicolumn{6}{l|}{\khoanh{$\ge 0$} / $\le 0$ / Tùy ý $\longrightarrow$ \khoanh{$\le$} / $\ge$ / $=$ } & \multicolumn{2}{c|}{ \multirow{4}{*}{\shortstack{\textit{Các cột này} \\ \textit{trở thành} \\ \textit{RBC đối ngẫu} \\ \textit{$R'_1, R'_2, R'_3$}}} } \\
        \cline{2-9}
        & \multicolumn{2}{c}{$x_2$} & \multicolumn{6}{l|}{\khoanh{$\ge 0$} / $\le 0$ / Tùy ý $\longrightarrow$ \khoanh{$\le$} / $\ge$ / $=$ } & \multicolumn{2}{c|}{} \\
        \cline{2-9}
        & \multicolumn{2}{c}{$x_3$} & \multicolumn{6}{l|}{\khoanh{$\ge 0$} / $\le 0$ / Tùy ý $\longrightarrow$ \khoanh{$\le$} / $\ge$ / $=$ } & \multicolumn{2}{c|}{} \\
        \hline
    \end{tabular}
    }
    \caption{Phân rã dữ kiện và cấu trúc dấu của bài toán Gốc (MIN)}
\end{table}

\begin{table}[H]
    \centering
    \renewcommand{\arraystretch}{1.8}
    \resizebox{\textwidth}{!}{
    \begin{tabular}{|c|c c c c c c c c|c|c|}
        \hline
        \textbf{Bài toán ĐỐI NGẪU} & \multicolumn{8}{c|}{\textbf{Hệ số \& Biến số ĐN ($y_i$)}} & \textbf{Dấu RBC ĐN} & \textbf{Hệ số $c$} \\
        \hline
        $Z' (\max)$ & \TL{b_1}{12} & $y_1$ & $+$ & \TL{b_2}{6} & $y_2$ & $+$ & \TL{b_3}{24} & $y_3$ & \multicolumn{2}{c|}{\textit{(Từ hệ số mục tiêu gốc)}} \\
        \hline
        \textbf{Biến gốc $\rightarrow$ RBC ĐN} & \multicolumn{8}{c|}{\textit{Ma trận hệ số chuyển vị $A^T$}} & \textbf{Dấu biến gốc $\rightarrow$ Dấu RBC ĐN} & \\
        \hline
        $x_1 \longrightarrow R'_1$ & \TL{a_{11}}{2} & $y_1$ & $+$ & \TL{a_{21}}{1} & $y_2$ & $+$ & \TL{a_{31}}{6} & $y_3$ & 
        \khoanh{$\ge 0$} / $\le 0$ / Tùy ý $\longrightarrow$ \khoanh{$\le$} / $\ge$ / $=$ & \TL{c_1}{12} \\
        \hline
        $x_2 \longrightarrow R'_2$ & \TL{a_{12}}{3} & $y_1$ & $+$ & \TL{a_{22}}{3} & $y_2$ & $+$ & \TL{a_{32}}{9} & $y_3$ & 
        \khoanh{$\ge 0$} / $\le 0$ / Tùy ý $\longrightarrow$ \khoanh{$\le$} / $\ge$ / $=$ & \TL{c_2}{27} \\
        \hline
        $x_3 \longrightarrow R'_3$ & \TL{a_{13}}{2} & $y_1$ & $+$ & \TL{a_{23}}{1} & $y_2$ & $+$ & \TL{a_{33}}{2} & $y_3$ & 
        \khoanh{$\ge 0$} / $\le 0$ / Tùy ý $\longrightarrow$ \khoanh{$\le$} / $\ge$ / $=$ & \TL{c_3}{6} \\
        \hline
        \multirow{4}{*}{\shortstack{\textbf{Dấu RBC gốc} \\ $\downarrow$ \\ \textbf{Dấu Biến ĐN}}} 
        & \multicolumn{2}{c}{$y_1$} & \multicolumn{6}{l|}{\khoanh{$\ge$} / $\le$ / $=$ $\longrightarrow$ \khoanh{$\ge 0$} / $\le 0$ / Tùy ý } & \multicolumn{2}{c|}{ \multirow{4}{*}{\shortstack{\textit{Các cột này} \\ \textit{thể hiện} \\ \textit{chiều biến} \\ \textit{đối ngẫu}}} } \\
        \cline{2-9}
        & \multicolumn{2}{c}{$y_2$} & \multicolumn{6}{l|}{\khoanh{$\ge$} / $\le$ / $=$ $\longrightarrow$ \khoanh{$\ge 0$} / $\le 0$ / Tùy ý } & \multicolumn{2}{c|}{} \\
        \cline{2-9}
        & \multicolumn{2}{c}{$y_3$} & \multicolumn{6}{l|}{\khoanh{$\ge$} / $\le$ / $=$ $\longrightarrow$ \khoanh{$\ge 0$} / $\le 0$ / Tùy ý } & \multicolumn{2}{c|}{} \\
        \hline
    \end{tabular}
    }
    \caption{Ánh xạ hệ số và thiết lập bài toán Đối ngẫu (MAX)}
\end{table}

\noindent Từ Bảng ánh xạ, ta dễ dàng viết được bài toán Đối ngẫu (D) một cách trực quan và chính xác:
\begin{align*}
    Z' &= 12y_1 + 6y_2 + 24y_3 \rightarrow \max \\
\text{sao cho:} \quad
    &2y_1 + y_2 + 6y_3 \le 12 \\
    &3y_1 + 3y_2 + 9y_3 \le 27 \\
    &2y_1 + y_2 + 2y_3 \le 6 \\
    &y_1, y_2, y_3 \ge 0
\end{align*}

\textbf{Cách sử dụng:}
\begin{enumerate}
    \item Đọc đề bài, xác định bài toán là MAX hay MIN để chọn đúng 1 trong 2 bảng trên.
    \item Điền các hệ số $a_{ij}$, $c_j$ và $b_i$ vào các vị trí tương ứng.
    \item Tại cột \textbf{Ánh xạ Dấu}, nhìn vào dấu của đề bài và dùng lệnh \verb|\khoanh{}| để đóng khung cụm chuyển đổi đúng, dùng \verb|\cancel{}| để gạch bỏ các cụm dư thừa.
    \item Đọc bảng theo chiều dọc để viết ra hệ Ràng buộc chính đối ngẫu ($R'_1, R'_2, R'_3$), và dọc theo hệ số $b$ để viết hàm mục tiêu đối ngẫu.
\end{enumerate}